%!TEX root = ../thesis.tex
\setlength{\parskip}{0pt}
\chapter{哈希检索方法研究现状}
本章围绕哈希检索方法研究现状展开综述。首先，介绍单模态与跨模态哈希检索的基本原理，给出哈希映射函数的定义、汉明距离的度量方式以及跨模态联合优化模型的基本数学形式，为后续方法分析与模型构建提供理论基础。然后，分别梳理静态与动态单模态哈希方法的代表性研究，重点分析其在概念漂移条件下的适应性与更新代价问题；最后，归纳静态与动态跨模态哈希方法的发展情况，讨论在线更新中的语义保持、跨模态对齐及历史信息保留等关键问题。基于上述分析，进一步形成第三章 SCPH 方法与第四章 MIH 方法的研究动机。
\section{哈希算法基础}

单模态哈希检索与跨模态哈希检索在方法本质上具有一致性：二者均通过学习哈希函数将原始高维特征映射为紧凑二值码，并在汉明空间中以近似最近邻方式实现高效检索；其核心目标均可归结为在低存储与低计算代价下保持语义相似性。然而，二者在问题设定与优化重心上存在显著差异。单模态哈希主要面向同构特征空间，重点在于保持同一模态内部的局部邻域结构与语义判别性；跨模态哈希则面向异构特征空间，除需保持模态内语义结构外，还需显式学习不同模态之间的语义对齐关系，以缩小“语义鸿沟”。因此，跨模态哈希通常引入共享子空间学习、跨模态一致性约束或联合表示学习等机制，其优化难度与建模复杂度相对更高。基于上述联系与差异，本文将本节内容划分为“单模态场景下的哈希检索基本原理”与“跨模态场景下的哈希检索基本原理”两个部分，分别展开阐述。

\subsection{单模态场景下的哈希检索基本原理}
哈希检索是一种将高维特征映射为低维二值编码的技术，其核心思想是通过设计哈希函数将高维特征空间中的样本映射到汉明空间，并基于汉明距离实现近似最近邻搜索。哈希检索的基本过程可表述为：对输入样本 $\mathbf{x}\in\mathbb{R}^{d}$ 依次施加 $K$ 个哈希函数，生成长度为 $K$ 的二值编码。将该过程表示为如下公式：
\begin{equation}
  \mathbf{H}(\mathbf{x})=[h_1(\mathbf{x}),h_2(\mathbf{x}),\dots,h_K(\mathbf{x})], \quad \mathbf{H}(\mathbf{x})\in\{-1,+1\}^{K}
\end{equation}
其中，$h_k(\mathbf{x})$ 表示第 $k$ 个哈希函数。由此，原始高维特征被映射到低维汉明空间，检索阶段可基于二值码之间的汉明距离实现近似最近邻搜索。由于编码紧凑且距离计算可由位运算高效完成，哈希检索在大规模场景中兼具较低存储开销与较高检索效率。

从函数构造方式看，哈希映射可采用线性形式、核映射形式以及深度网络形式。其中，线性哈希函数因形式简单、训练与推理开销低而被广泛使用，具体公式如下：
\begin{equation}
  \mathbf{b}=\operatorname{sgn}(\mathbf{W}^{\top}\mathbf{x}), \quad \mathbf{b}\in\{-1,+1\}^{K}
\end{equation}
其中，$\mathbf{W}\in\mathbb{R}^{d\times K}$ 为投影矩阵。对于任意样本 $\mathbf{x}_i$ 与 $\mathbf{x}_j$，可通过公式(2-2)得到其对应的二值码 $\mathbf{b}_i,\mathbf{b}_j$，二者之间的汉明距离计算公式如下：
\begin{equation}
  D_H(\mathbf{b}_i,\mathbf{b}_j)=\frac{1}{2}\left(K-\mathbf{b}_i^{\top}\mathbf{b}_j\right)
\end{equation}
一般情况下，语义相近的样本之间的汉明距离较小，语义差异较大的样本之间的汉明距离较大。
\subsection{跨模态场景下的哈希检索基本原理}

跨模态哈希检索旨在将图像、文本等异构模态数据映射到统一汉明空间，使语义一致的跨模态样本在二值码空间中保持邻近关系。设图像模态样本集为 $\mathbf{X}=\{\mathbf{x}_i\}_{i=1}^{n}$、文本模态样本集为 $\mathbf{Y}=\{\mathbf{y}_i\}_{i=1}^{n}$，哈希码长度记为 $K$。跨模态哈希通过学习两类模态相关映射函数，实现从异构特征空间到共享二值空间的投影：
\begin{equation}
  \mathbf{b}_i^{x}=\operatorname{sgn}\!\left(f_x(\mathbf{x}_i;\theta_x)\right),\quad
  \mathbf{b}_i^{y}=\operatorname{sgn}\!\left(f_y(\mathbf{y}_i;\theta_y)\right),\quad
  \mathbf{b}_i^{x},\mathbf{b}_i^{y}\in\{-1,+1\}^{K}
\end{equation}
其中，$f_x(\cdot)$ 与 $f_y(\cdot)$ 分别表示图像与文本的哈希映射函数，$\theta_x,\theta_y$ 为对应模型参数。该过程的关键不在于仅压缩特征维度，而在于通过共享语义约束缩小模态间表示差异。

与单模态检索类似，跨模态检索仍以汉明距离作为相似性度量。对于任意查询码 $\mathbf{b}_q$ 与候选码 $\mathbf{b}_j$，其汉明距离可表示为：
\begin{equation}
  D_H(\mathbf{b}_q,\mathbf{b}_j)=\frac{1}{2}\left(K-\mathbf{b}_q^{\top}\mathbf{b}_j\right)
\end{equation}
因此，跨模态哈希的优化目标可归结为：使语义相似的跨模态样本对具有更小的汉明距离，语义不相似样本对具有更大的汉明距离。

为同时保证语义保持、跨模态对齐与二值化质量，常用联合优化框架可表示为如下形式\upcite{Ding2014CMFH,Lin2015SePH,Zhang2014SCM,Jiang2017DCMH}：
\begin{equation}
  \begin{aligned}
    \min_{\theta_x,\theta_y,\mathbf{B}} \quad
    & \mathcal{L}_{sem}(\mathbf{U},\mathbf{V},\mathbf{S})
    +\lambda_1\mathcal{L}_{align}(\mathbf{U},\mathbf{V}) \\
    & +\lambda_2\!\left(\|\mathbf{U}-\mathbf{B}\|_F^2+\|\mathbf{V}-\mathbf{B}\|_F^2\right)
    +\lambda_3\mathcal{R}(\mathbf{B}) \\
    \text{s.t.} \quad
    & \mathbf{B}\in\{-1,+1\}^{n\times K}
  \end{aligned}
\end{equation}
其中，$\mathbf{U}=f_x(\mathbf{X};\theta_x)$、$\mathbf{V}=f_y(\mathbf{Y};\theta_y)$ 为连续表示，$\mathbf{S}$ 为跨模态语义关系矩阵；$\mathcal{L}_{sem}$ 用于保持语义相似结构，$\mathcal{L}_{align}$ 用于约束跨模态表示一致性，量化项用于减小连续表示与二值码之间的误差，$\mathcal{R}(\mathbf{B})$ 通常包含比特平衡与去相关约束。

\section{单模态哈希检索算法}

\subsection{静态单模态哈希方法}
静态哈希算法根据模型训练是否依赖数据可分为数据独立哈希和数据依赖哈希。数据独立哈希（Data Independent Hashing）方法的核心特征是哈希函数的设计不依赖于数据集的具体分布。换句话说，这类方法假设数据的分布是未知的，哈希函数的构造仅依赖于预设的数学原理或随机化过程，而不是数据本身。最具代表性的数据独立哈希方法是局部敏感哈希LSH\upcite{LSH}, 其基本思想是通过构造一组哈希函数，使得相似的数据点在哈希空间中具有较高的碰撞概率，而不相似的点则具有较低的碰撞概率。LSH方法的主要优势在于其理论可解释性和应用广泛性，尤其适用于大规模的近似最近邻搜索（ANN）。但是由于LSH中采用了随机映射算法，因此索引结果不可控。初始的LSH是用来解决汉明空间的高维搜索问题，基于p-stable的LSH\upcite{datar2004locality}改进了LSH算法，将LSH的使用范围扩展到了欧氏空间。尽管数据独立哈希方法在计算和存储上具有较高的效率，但它们的缺点也十分明显。由于不考虑数据的具体分布，数据独立哈希方法通常无法充分挖掘数据的内在结构，导致其在某些复杂数据集上的效果不如预期。此外，数据独立哈希方法往往无法很好地应对数据分布的动态变化，尤其是在动态数据环境中，其性能可能会显著下降。

相比于数据独立哈希，数据依赖哈希（Data Dependent Hashing）方法则通过学习数据的分布结构来构建哈希函数，从而生成能够更好地表示数据特征的哈希码。数据依赖哈希方法\upcite{weiss2008spectral, xu2011complementary, strecha2011ldahash, gong2012iterative, shen2015supervised, gui2017fast, gui2016supervised}不仅依赖于数据集的具体内容，还能够在训练过程中优化哈希函数，以增强相似样本之间的区分度。这些方法通常利用无监督学习、监督学习或半监督学习等技术，从数据中学习到能够保留数据在原始特征科技中的相似关系的哈希函数。

无监督哈希算法在训练时不依赖数据的标签信息，仅仅利用数据的特征信息来学习哈希函数。谱哈希（Spectral Hashing）\upcite{weiss2008spectral} 是无监督哈希算法中的典型方法之一，谱哈希通过构建是数据点之间的相似性图进行谱分解来学习数据的哈希编码。迭代量化哈希（ITQ）\upcite{gong2012iterative}将数据投影到一个超立方体上，学习一个最优的旋转矩阵，使得数据通过旋转矩阵能够映射到超立方体的顶点上，保留数据的结构信息。ITQ加强了哈希映射的均衡性，减小了量化误差，提高了检索精确度。但迭代量化的复杂度较高，同时由于迭代量化哈希算法需要先将数据投影到子空间后再减小量化误差，导致其编码位数不能超过样本的维数。
有监督哈希算法利用数据的标签语义信息和数据的特征信息共同学习哈希函数，数据的标签带有丰富的语义信息，能够有效地应对“语义鸿沟”问题，更好地实现哈希编码对图像语义信息的保留。相比于无监督哈希算法，有监督算法通常具有更优秀的检索性能。LDA\upcite{strecha2011ldahash} 哈希算法的目标是通过最小化哈希映射后相似样本之间的差异，同时最大化不相似样本之间的差异。在求解这个目标函数时，LDA哈希借鉴了线性判别分析方法，通过优化过程使得在哈希空间中，相似样本的哈希编码尽量接近，而不相似样本的哈希编码尽量远离，从而提高哈希编码的区分能力。有监督离散哈希算法（SDH）\upcite{shen2015supervised} 通过最小化分类器的损失函数来生成图像的哈希编码，以确保相同类别的样本在哈希空间中的编码尽可能相似，而不同类别的样本则尽量远离。为了进一步提高SDH的效果，SDHR\upcite{gui2016supervised}算法通过直接根据数据学习回归目标来优化哈希编码，使得哈希映射更加精确和高效。基于SDH的框架，研究人员还提出了一种快速SDH算法\upcite{gui2017fast}，旨在避免传统的迭代训练过程，从而加速哈希编码的学习速度，提升大规模数据集上的学习效率。

尽管有监督哈希算法能够拥有令人满意的检索效果，但是这类方法要求训练数据带有标签，而现实世界中要求所有的数据都带有标签往往是不现实的，尤其是在海量数据的情况下。半监督哈希算法结合了有监督哈希和无监督哈希的优点，能够利用少部分带有标签信息的数据和大部分无标签的数据训练得到一个高效的哈希模型，适应真实世界的数据场景。半监督哈希（SSH）\upcite{wang2012semi}提出了三种不同的半监督哈希框架，为后续半监督哈希算法的研究提供了有价值的参考。SSH\upcite{wang2012semi}算法逐个的训练哈希函数，每个哈希函数的训练都会侧重于改善前一个哈希函数没有训练好的样本相似度关系。基于 SSH 算法思想，BSPLH\upcite{wu2012semi}算法逐个训练哈希函数，每个哈希函数是根据之前所有哈希函数没有训练好的样本相似关系进行训练。

\subsection{动态单模态哈希方法}
在实际应用中，数据通常以数据流或连续数据块的形式持续产生。在这种动态数据环境下，传统的哈希方法往往仅依赖初始数据进行训练，难以应对数据的持续变化和更新。因此，动态哈希方法成为解决这一问题的有效策略。动态哈希方法能够在数据不断变化的过程中，实时更新哈希编码和结构，以适应新数据的加入，从而提高检索效率和系统的灵活性。

目前已有一些动态哈希算法提出，OKH\upcite{huang2017online}是第一个用于解决动态环境下的有监督哈希算法，其通过成对出现的新样本和它们之间的标签信息动态更新哈希函数。MIHash\upcite{cakir2017mihash}提出基于互信息的在线哈希学习，通过最大化近邻与非近邻样本在汉明空间中的可分性来更新哈希函数，并仅在检索质量提升时更新哈希表，从而减少重编码开销并兼顾检索精度与效率。无监督的OSH\upcite{leng2015online}算法提出数据在线抽象(Sketch)的思想，将原始数据在线更新为更低维度的数据抽象，保留原始数据的语义信息，然后进行高效的在线哈希函数学习。FROSH\upcite{chen2019making}在OSH的基础上，进一步优化，达到了更高、更快的检索精度和更低的存储开销，并提出了分布式的FROSH方法。有监督的HMOH算法\upcite{lin2020hadamard}使用启发式生成的Hadamard矩阵，根据新数据样本的标签信息为样本分配目标编码，指导哈希函数在线高效的更新。在线比特选择算法（Online Bit Selection Hashing）\upcite{weng2020online}根据样本的标签随机生成目标编码，用于训练R个哈希函数，使用比特选择算法选择最优r位（r << R）哈希函数构成哈希投影矩阵。无监督的在线球面哈希(Online Spherical Hashing)\upcite{weng2019fast}提出了一种基于球面几何的在线哈希算法，利用球面距离来定义数据样本间的相似度，这种方法不仅保留了传统哈希方法的低维二进制表示特点，还利用球面几何提供的优越性，在保证高效存储和检索的同时，增强了相似样本之间的区分能力。

上述的动态哈希方法能够在线更新哈希函数，以适应不断变化的数据流，但它们通常忽略了概念漂移（Concept Drift）现象\upcite{elwell2011incremental, ng2016incremental}。动态数据环境下，概念漂移现象不可避免的会发生，如新的数据类别出现、已有的数据类别分布发生改变。在动态哈希方法中，更新哈希函数虽然能够适应新样本的加入，但往往无法及时捕捉到数据分布的根本变化，尤其是在数据特征发生显著变化时。这种忽视概念漂移的情况可能导致哈希函数的失效，使得模型无法准确地反映当前数据的特征分布，从而影响检索的效果和准确性。ICH\upcite{ng2016incremental}算法通过多重哈希机制保留随着时间推移到达的新图像数据中的知识，并采用基于权重的排名方法，使得检索结果能够适应不断变化的数据环境。ICH是处理概念随时间漂移的开创性哈希方法之一，并获得了良好的性能。IBL\upcite{ng2018incremental}算法则通过在概念漂移发生时，利用三分量目标函数的迭代最大化过程，选择现有和新训练的哈希位，并为其加权，从而提升语义图像检索的准确性。互补增量哈希算法（CIHR）\upcite{tian2020complementary}通过增量方式训练多个哈希表，以应对动态数据的变化。CPH\upcite{tian2019concept}算法则不需要更新旧的哈希码，而是将来自不同时间或分布的相同概念图像映射到哈希投影空间的相似区域，从而有效处理跨时间段或分布的图像检索任务。无监督的动态UMH\upcite{yang2023unsupervised}哈希算法基于ITQ算法训练一个多哈希表系统，保留新旧数据的知识知识。DIH\upcite{tian2023deep}在深度网络中使用了一种多哈希表框架来同时保留并学习以前的知识和新的数据分布信息，实现了对概念漂移数据环境的适应。这些方法均致力于解决动态数据环境下哈希检索中的挑战，尤其是概念漂移带来的问题，提高了图像检索的灵活性和准确性。

总体而言，现有方法在概念漂移场景下仍存在两点不足：一是多数方法在模型更新后仍需访问历史数据并进行大规模重新编码，系统开销较高；二是部分方法对全监督标签依赖较强，在真实数据流中难以满足。针对上述问题，本文第三章提出SCPH方法：在保留概念信息的同时引入半监督学习，联合利用少量标注样本与大量未标注样本进行模型更新，并以类别概念编码进行保持与迁移，从而在降低人工标注依赖的同时缓解重编码负担，提升概念漂移环境下单模态哈希检索的实用性与鲁棒性。

\section{跨模态哈希检索算法}
\subsection{静态跨模态哈希方法}
静态跨模态哈希方法通常建立在数据环境稳定不变的假设之上，并利用全部训练数据进行离线建模与哈希函数学习。依据是否使用标签语义信息，现有离线跨模态哈希方法可进一步划分为无监督跨模态哈希与有监督跨模态哈希两类。

无监督跨模态哈希方法通常通过将多模态特征映射到共享汉明空间来学习哈希函数\upcite{Ding2014CMFH, wang2017robust, tu2023unsupervised, lu2023deep}。其中，CMFH\upcite{Ding2014CMFH}是较早实现不同模态统一哈希码与哈希函数联合学习的代表性方法，其核心思想是利用协同矩阵分解进行共享表示建模。随后，相关研究进一步引入离散优化与结构保持机制：如 RFDH\upcite{wang2017robust} 通过离散协同分解生成二值码，并结合分类框架学习哈希映射；UCHM\upcite{tu2023unsupervised} 则在哈希相似性保持损失下引入模态交互式相似性生成器，并结合比特选择模块挖掘跨模态编码交互信息，从而提升无监督跨模态检索效果。

相较而言，有监督跨模态哈希方法利用标签语义信息进行哈希学习，通常能够获得更优检索性能\upcite{hu2024cross, wu2024contrastive, lan2022label, Wang2021BATCH, Wang2022FCMH}。例如，SCH\upcite{hu2024cross} 将样本对划分为完全语义正相关、部分语义正相关与语义负相关等类型，并施加自适应约束以提高汉明空间利用效率；LGDH\upcite{lan2022label} 融合局部类别分布流形结构与多标签平衡矩阵，在不依赖连续松弛的条件下直接优化二值码；FCMH\upcite{Wang2022FCMH} 则联合全局与局部相似性嵌入，通过离散优化直接学习跨模态二值表示，在检索精度与训练效率之间取得较好平衡。

\subsection{动态跨模态哈希方法}
在静态数据环境中，跨模态哈希方法通常基于全量样本进行离线训练，能够在既定数据分布下取得较好的检索效果。然而，在实际应用场景中，数据往往呈现随时间不断变化的动态特性。当数据分布发生演化时，传统静态方法由于缺乏增量更新机制，难以及时对模型参数进行调整，导致已学习的哈希表示逐渐偏离当前数据分布，从而引起检索性能下降。为适应动态数据环境下的持续更新需求，在线跨模态哈希方法逐渐成为研究关注的重点方向。

在无监督在线跨模态哈希方面，代表性方法主要通过跨模态共享表示的增量学习来提升模型更新效率。OCMH\upcite{xie2016online} 通过学习共享潜在编码（Shared Latent Code, SLC），在不断到达的新数据上实现哈希码的高效更新，从而兼顾检索性能与在线学习效率。OCMFH\upcite{wang2020online} 则可视为 CMFH\upcite{xie2016online} 在动态环境下的在线扩展，通过持续学习流式样本信息，缓解离线方法在动态场景中的性能退化问题。总体而言，这类方法在缺乏标签监督的条件下具备较好的实用性，但受限于语义约束不足，其检索精度通常仍有提升空间。

相较之下，监督在线跨模态哈希方法通过引入标签信息强化语义保持能力，整体性能通常优于无监督方法\upcite{yao2019online, yang2023efficient, li2024robust, kang2023online, shu2023online, peng2024olch, kang2024discrete}。例如，OLSH\upcite{yao2019online} 将标签信息嵌入连续潜在语义空间，并基于新到达样本动态学习哈希函数；OSCMFH\upcite{yang2023efficient} 利用标签嵌入构建哈希码，有效融合了跨模态共享语义与模态特有语义表示。进一步地，ROHLSE\upcite{li2024robust} 通过低秩与稀疏约束挖掘样本与标签之间的内在依赖关系，实现噪声标签纠正与标签预测；SHOH\upcite{han2024supervised} 则结合层次化标签相似性与加权检索策略，提升语义表征能力，并缓解流式跨模态检索中的比特翻转误差。由此可见，如何在在线更新效率与语义判别能力之间取得平衡，仍是监督在线跨模态哈希研究中的关键问题。

尽管在线跨模态哈希方法已有较多研究，但现有工作主要关注利用新到达样本对模型进行增量更新，对动态数据环境中普遍存在的概念漂移问题考虑不足。概念漂移发生后，数据分布随时间变化，现有方法往往难以及时调整表示与映射关系，进而造成检索性能下降。另一方面，多数方法基于单哈希表进行持续更新，容易产生灾难性遗忘，导致历史语义信息保留不足。针对上述问题，本文在第四章提出基于多表的增量式哈希方法（Multi-table based Incremental Hashing, MIH），以增强模型在动态环境下对分布变化的适应能力，并提高跨模态检索的整体性能。
